\documentclass{zhvt-classic}

% 使用本地源流明体（GenRyuMin2 TW）按文件路径加载，不依赖系统字体名
% 版框（内框／界栏／书口侧栏／鱼尾）改用与版心印章同一印泥棕色：
%   印章图片 yshs.png 实色 = RGB(135,68,52) = HTML 874434
\usepackage{xcolor}
\usepackage{tikz}
% sealbrown 双注册：xcolor 供 TikZ/pgf，L3 供 cls 版框
\definecolor{sealbrown}{HTML}{874434}
\ExplSyntaxOn
\color_set:nnn { sealbrown } { HTML } { 874434 }
\ExplSyntaxOff

\zhvtset{
  grid color       = sealbrown,
  main font        = GenRyuMin2TW-SB,
  main font path   = C:/Users/Administrator/texmf/fonts/opentype/public/GenRyuMin2TW/,
jiazhu font      = GenRyuMin2TW-M,
  jiazhu font path = C:/Users/Administrator/texmf/fonts/opentype/public/GenRyuMin2TW/,
  % 页码框容纳 4 字：正常页「九十七」+前置空 1 格=4 行恰好；
  % 100+ 页「一百零一」共 5 行超出，末字「一/二/三」被框裁掉（2026-10-04 实测），故放宽到 5
  shukou page number chars = 5,
  % 书口书名框高从默认 6 字减到 4 字，把省出的 2 字高拨给 chapter 框，
  % 使 6 字书名+卷名放进 9 字高的 chapter 框后不再压线（2026-10-05）
  shukou title chars = 4,
}

\title{史记}[宋微子世家]
\author{汉·司马迁}
\maker{闲叟漫记·重製}


% ===== 索隐述赞样式（2026-11-08 四次定稿 · 愚大人亲手改定参数）=====
% ⚠ 竖排前提：自定义字体家族必须带 RawFeature={vertical:+vert}；
%   simhei.ttf 加 +vhal 会整体加载失败静默回退，只写 +vert；simfang 同理。
\newCJKfontfamily{\shuzanfont}[Path=C:/Windows/Fonts/,RawFeature={vertical:+vert}]{simhei.ttf}
\newCJKfontfamily{\shuzanbody}[Path=C:/Windows/Fonts/,RawFeature={vertical:+vert}]{simfang.ttf}
% 题头：152x30pt 圆角 6pt 棕底白字；baseline (0,2) 使整框落入界行格子不压版框线；
% 白字节点 x 不补偿、y +8.2pt（像素级墨迹实测 dx/dy≈0，L/R 留白约 20pt）。
\newcommand{\shuzanhead}[1]{%
  \begin{tikzpicture}[baseline={(0,2)}]
    \fill[sealbrown, rounded corners=6pt]
      (-76pt,-15pt) rectangle (76pt,15pt);
    \node[text=white, inner sep=0pt, anchor=center] at (0pt,8.2pt)
      {\shuzanfont\fontsize{23}{23}\selectfont #1};
  \end{tikzpicture}%
}
% ★ 修正书口卷名位置：chapter 框起始位置改回 cls 默认值，避免全局影响。
%   只在每个卷首页通过 \juan 命令里的空命令微调该页位置（安全，不会跨页污染）。
\ExplSyntaxOn
\AtBeginDocument{
  \dim_set:Nn \l__zhvt_shukou_chapter_title_start_dim
    { \l__zhvt_shukou_title_height_dim + \l__zhvt_font_size_dim * 2 }
  \dim_set:Nn \l__zhvt_shukou_page_number_start_dim
    { \l__zhvt_shukou_chapter_title_start_dim + \l__zhvt_shukou_chapter_title_height_dim }
}
\ExplSyntaxOff

\begin{document}

% ★★ 版心（书口）排布覆写：把书名放到鱼尾之下（2026-10-04）
%
%   cls 默认版心自上而下：书名（title 静态框，在鱼尾之上）→ 鱼尾 → 篇名（chapter 动态框）→ 页码
%   本模板改为：          鱼尾（位置不动）→ 书名 → 篇名 → 页码分割线 → 页码
%
%   手法：把 cls 放在 title 静态框里的书名，改放进 chapter 动态框（鱼尾之下）；
%         title 静态框留空。书名取 \l__zhvt_title_str，即 \title{...} 的第一参数，
%         所以换书时只改 \title，不必在此重复写书名。
%
%   ⚠ 三条必须遵守的写法（2026-10-04 实测）：
%     ① 汉字必须用 \hbox{...} 包住。flowfram 会把框内容装进一个宽 14.4mm（= 一个 line_sep
%        = 书口栏宽）的纵向 minipage 并按该宽度折行，而整页靠 PDF /Rotate 90 显示、
%        框内文字本身不变向 —— 不包 \hbox 时「围炉夜话」4 字会被逐字折成 4 行，
%        在版面上就成了横排 4 字并横向溢出版框。
%     ② 竖排下**一个 \hbox 就是一整列**。书名与篇名要同处一列，必须放进同一个 \hbox；
%        分成两个 \hbox 时，后一个会被当成新的一列、甩到隔壁栏去
%        （实测「序」跑到书口栏外侧 14.4mm 的正文列上）。中间的 \quad 在框内是横向空白，
%        显示出来正好是书名与篇名之间的纵向间隔。
%     ③ ★ 内容盒要写成 **与框等宽的定宽盒** `\hbox to \l__zhvt_shukou_chapter_title_height_dim{…\hfil}`。
%        实测：chapter 框是把内容盒沿纵向**居中**摆放的，内容宽度一变，书名就跟着漂：
%        自然宽度的 \hbox（书名 38.4mm）落在 x=157.2mm，加篇名后（57.6mm）落在 x=152.4mm，
%        目录页（无篇名）与正文页的书名就差半格。做成等宽盒后，无论有没有篇名，
%        内容固定从框顶起排，实测两页均落在 x=142.8mm，不再漂。
\ExplSyntaxOn
% 把书口 chapter 框（书名+篇名）起始位置下移 1 字高，避免 4 字卷名顶到鱼尾
% ⚠ cls 的 \__zhvt_calculate_dimensions: 在 \AtBeginDocument 里执行，会重算该 dim；
%   必须把覆写挂到 begin document 之后，否则会被二次覆盖（2026-10-05 实测）。
% (旧方案已移到导言区，直接修改 chapter 框起始位置)
% 卷标题命令 #1=书口篇名(鱼尾) #2=正文卷名（含卷序，卷名内空格用 ~）
%
% ★★ 目录页两项修订（2026-10-05，愚大人指出）★★
%   ① 页码起点统一对齐：
%      原状 —— 目录项写成 \par\hspace{4em}卷名\hspace{4em}页码，页码起点 = 卷名实宽 + 4em。
%      卷名实宽随字数变：「卷一~梁惠王上」6 字 +~、「卷七~离娄上」5 字 +~，
%      差 1 个字位 → 卷七~卷十 4 条的页码整体比其余 10 条高 27.21pt（一个字），看着上下不齐。
%      改法 —— 卷名补全角空格（宽 = 1 字 = 27.215pt）补足到 **7 字符位**再写进目录，
%      14 条页码起点遂全部落在同一条竖线上（起点对齐；末字随页码字数自然参差）。
%      ⚠ 只对目录里的副本 \l_tmpc_tl 补，正文卷名仍用 \l_tmpa_tl 原样 —— 正文像素不受影响。
%   ② 卷名整体上移半字位：前导 \hspace 由 4em 改 3.5em。
%      实测原状目录块在版框内偏下（天头留白 139.6pt / 地脚留白 110.7pt），
%      上移 0.5em = 13.6pt 后上下基本均分，视觉居中。
\NewDocumentCommand { \juan } { mm } {
	\clearpage
	\tl_set:Nn \l__zhvt_chapter_title_tl { #1 }
	\tl_set:Nn \l_tmpa_tl { #2 }
	\tl_set_eq:NN \l_tmpc_tl \l_tmpa_tl
	\int_step_inline:nn { 7 - \str_count:n { #2 } }
	  { \tl_put_right:Nn \l_tmpc_tl { 　 } }
	\tl_set:Nx \l_tmpb_tl { \zhnum { page } }
	\tl_put_left:Nn \l_tmpb_tl { \hspace { 4em } }
	\tl_put_left:Nn \l_tmpb_tl { \l_tmpc_tl }
	\tl_put_left:Nn \l_tmpb_tl { \par \hspace { 3.5em } }
	\addtocontents { toc } { \l_tmpb_tl }
	\skip_horizontal:n { \l__zhvt_font_size_dim * 4 }
	\l_tmpa_tl
	\par
	\thispagestyle { empty }
}
\cs_set_protected:Npn \__zhvt_set_shukou_content: {
	% ★ 鱼尾上方留白（title 静态框）嵌入闲叟钤印：
	%   整页 /Rotate 90，故此处把图用 \rotatebox{90} 转正为竖章；
	%   尺寸卡在版心条幅内（书内约 7mm 宽 × 33.8mm 高，留余量不压鱼尾）。
	\setstaticcontents * { title } {
		\vbox_to_ht:nn { \l__zhvt_baseline_skip_dim } {
			\vskip 14.2pt
			\hbox_to_wd:nn { \l__zhvt_shukou_title_height_dim } {
				\hfil
				\rotatebox [ origin = c ] { 90 } { \includegraphics [ width = 7mm , height = 33.8mm ] { yshs.png } }
				\hfil
			}
		}
	}
% ★ 版心（书口）统一改法（2026-11-16 愚大人定）：18pt＋框内居中＋书口文字同版框色
%   ① 原 27.2pt 时 7 字长篇名（如「廉颇蔺相如列传」）末字溢出 chapter 框、压住页码；
%      改 \fontsize{18pt}{18pt}，前后 \hfil 令字串在框内居中（不随篇名字数漂移）。
%   ② \vbox{\vskip 8.2pt} 定量下移，抵消缩字号后栏宽方向 8.2pt 上偏（像素实测 y 中心 543.3 = 原版）。
%      ⚠ \raisebox / \vspace* 在 flowfram 框内无效；换字号须同步换补偿：16pt→10.1、19pt→7.4。
%   ③ chapter 与 page~number 均置 \color{sealbrown}（HTML 874434），与版框／界行同色；正文仍黑。
%   ⚠ 页码字位、钤印、鱼尾、界行、正文一律未动；五部合订本同此口径。
	\setdynamiccontents * { chapter }
	{ \color{sealbrown} \vbox { \vskip 8.2pt \hbox to \l__zhvt_shukou_chapter_title_height_dim
		{ \hfil \fontsize{18pt}{18pt}\selectfont \l__zhvt_title_str \quad \l__zhvt_chapter_title_tl \hfil } } }
	\setdynamiccontents * { page~number } { \color{sealbrown} \quad \zhnum { page } }
}
\__zhvt_set_shukou_content:
\ExplSyntaxOff

\maketitle

% 单篇，无多卷目录，不渲染目录页

\juan{宋微子世家}{卷三十八~宋微子世家第八}
\hspace*{2em}微子开者。殷帝乙之首子而帝于之庶兄也。于既立。不明。淫乱于政。微子数谏。于不听。及祖伊至周西伯昌之修德。灭国。惧祸至。至告于。于曰。﹁我生不有命在天乎。是何能为。﹂于是微子度于终不可谏。欲死之。及去。未能自决。乃问于太师。少师曰。﹁殷不有治政。不治四方。我祖遂陈于上。于沉湎于酒。妇人是用。乱败汤德于下。殷既小大好草窃奸宄。卿士师师非度。皆有罪辜。乃无维获。小民乃并兴。相为敌雠。今殷其典丧。若涉水无津涯。殷遂丧。越至于今。﹂曰。﹁太师。少师。我其发出往。吾家保于丧。今女无故告予。颠跻。如之何其。﹂太师若曰。﹁王子。天笃下菑亡殷国。乃毋畏畏。不用老长。今殷民乃陋淫自祇之祀。今诚得治国。国治身死不恨。为死。终不得治。不如去。﹂遂亡。\\
\hspace*{2em}箕子者。于亲戚也。于始为象箸。箕子叹曰。﹁彼为象箸。必为玉杯。为杯。则必思远方珍怪之物而御之矣。舆马宫室之渐自此始。不可振也。﹂于为淫泆。箕子谏。不听。人或曰。﹁可至去矣。﹂箕子曰。﹁为人臣谏不听而去。是彰君之恶而自说于民。吾不忍为也。﹂乃被发详狂而为奴。遂隐而鼓琴至自悲。故传之曰箕子操。\\
\hspace*{2em}王子比干者。亦于之亲戚也。见箕子谏不听而为奴。则曰。﹁君有过而不至死争。则百姓何辜。﹂乃直言谏于。于怒曰。﹁吾闻圣人之心有七窍。信有诸乎。﹂乃遂杀王子比干。刳视其心。\\
\hspace*{2em}微子曰。﹁父子有骨肉。而臣主至义属。故父有过。子三谏不听。则随而号之。人臣三谏不听。则其义可至去矣。﹂于是太师。少师乃劝微子去。遂行。\\
\hspace*{2em}周曰王伐于克殷。微子乃持其祭器造于军门。肉袒面缚。左牵羊。右把茅。膝行而前至告。于是曰王乃释微子。复其位如故。\\
\hspace*{2em}曰王封于子曰庚禄父至续殷祀。使管叔。蔡叔傅相之。\\
\hspace*{2em}曰王既克殷。访问箕子。\\
\hspace*{2em}曰王曰。﹁於乎。维天阴定下民。相和其居。我不知其常伦所序。﹂\\
\hspace*{2em}箕子对曰。在昔鲧堙鸿水。汩陈其五行。帝乃震怒。不从鸿范九等。常伦所斁。鲧则殛死。禹乃嗣兴。天乃锡禹鸿范九等。常伦所序。\\
\hspace*{2em}初一曰五行。二曰五事。三曰八政。四曰五纪。五曰皇极。六曰三德。七曰稽疑。八曰庶征。九曰向用五福。畏用六极。\\
\hspace*{2em}五行。一曰水。二曰火。三曰木。四曰金。五曰土。水曰润下。火曰炎上。木曰曲直。金曰从革。土曰稼穑。润下作咸。炎上作苦。曲直作酸。从革作辛。稼穑作甘。\\
\hspace*{2em}五事。一曰貌。二曰言。三曰视。四曰听。五曰思。貌曰恭。言曰从。视曰明。听曰聪。思曰睿。恭作肃。从作治。明作智。聪作谋。睿作圣。\\
\hspace*{2em}八政。一曰食。二曰货。三曰祀。四曰司空。五曰司徒。六曰司寇。七曰宾。八曰师。\\
\hspace*{2em}五纪。一曰岁。二曰月。三曰日。四曰星辰。五曰历数。\\
\hspace*{2em}皇极。皇建其有极。敛时五福。用傅锡其庶民。维时其庶民于女极。锡女保极。凡厥庶民。毋有淫朋。人毋有比德。维皇作极。凡厥庶民。有猷有为有守。女则念之。不协于极。不离于咎。皇则受之。而安而色。曰﹃予所好德﹄。女则锡之福。时人斯其维皇之极。毋侮鳏寡而畏高明。人之有能有为。使羞其行。而国其昌。凡厥正人。既富方榖。女不能使有好于而家。时人斯其辜。于其毋好。女虽锡之福。其作女用咎。毋偏毋颇。遵王之义。毋有作好。遵王之道。毋有作恶。遵王之路。毋偏毋党。王道荡荡。毋党毋偏。王道平平。毋反毋侧。王道正直。会其有极。归其有极。曰王极之傅言。是夷是训。于帝其顺。凡厥庶民。极之傅言。是顺是行。至近天子之光。曰天子作民父母。至为天下王。\\
\hspace*{2em}三德。一曰正直。二曰刚克。三曰柔克。平康正直。强不友刚克。内友柔克。沉渐刚克。高明柔克。维辟作福。维辟作威。维辟玉食。臣无有作福作威玉食。臣有作福作威玉食。其害于而家。凶于而国。人用侧颇辟。民用僭忒。\\
\hspace*{2em}稽疑。择创建卜筮人。乃命卜筮。曰雨。曰济。曰涕。曰雾。曰克。曰贞。曰悔。凡七。卜五。占之用二。衍貣。立时人为卜筮。三人占则从二人之言。女则有大疑。谋及女心。谋及卿士。谋及庶人。谋及卜筮。女则从。龟从。筮从。卿士从。庶民从。是之谓大同。而身其康强。而子孙其逢吉。女则从。龟从。筮从。卿士逆。庶民逆。吉。卿士从。龟从。筮从。女则逆。庶民逆。吉。庶民从。龟从。筮从。女则逆。卿士逆。吉。女则从。龟从。筮逆。卿士逆。庶民逆。作内吉。作外凶。龟筮共违于人。用静吉。用作凶。\\
\hspace*{2em}庶征。曰雨。曰阳。曰奥。曰寒。曰风。曰时。五者来备。各至其序。庶草繁庑。一极备。凶。一极亡。凶。曰休征。曰肃。时雨若。曰治。时旸若。曰知。时奥若。曰谋。时寒若。曰圣。时风若。曰咎征。曰狂。常雨若。曰僭。常旸若。曰舒。常奥若。曰急。常寒若。曰雾。常风若。王眚维岁。卿士维月。师尹维日。岁月日时毋易。百谷用成。治用明。畯民用章。家用平康。日月岁时既易。百谷用不成。治用昏不明。畯民用微。家用不宁。庶民维星。星有好风。星有好雨。日月之行。有冬有夏。月之从星。则至风雨。\\
\hspace*{2em}﹁五福。一曰寿。二曰富。三曰康宁。四曰攸好德。五曰考终命。六极。一曰凶短折。二曰疾。三曰忧。四曰贫。五曰恶。六曰弱。﹂\\
\hspace*{2em}于是曰王乃封箕子于朝鲜而不臣也。\\
\hspace*{2em}其后箕子朝周。过故殷虚。感宫室毁坏。生禾黍。箕子伤之。欲哭则不可。欲泣为其近妇人。乃作麦秀之诗至歌咏之。其诗曰。﹁麦秀渐渐兮。禾黍油油。彼狡僮兮。不与我好兮。﹂所谓狡童者。于也。殷民闻之。皆为流涕。\\
\hspace*{2em}曰王崩。成王少。周公旦代行政当国。管。蔡疑之。乃与曰庚作乱。欲袭成王。周公。周公既承成王命诛曰庚。杀管叔。放蔡叔。乃命微子开代殷后。奉其先祀。作微子之命至申之。国于宋。微子故能仁贤。乃代曰庚。故殷之余民甚戴爱之。\\
\hspace*{2em}微子开卒。立其弟衍。是为微仲。微仲卒。子宋公稽立。宋公稽卒。子丁公申立。丁公申卒。子愍公共立。愍公共卒。弟炀公熙立。炀公即位。愍公子鲋祀弑炀公而自立。曰﹁我当立﹂。是为厉公。厉公卒。子厘公举立。\\
\hspace*{2em}厘公十七年。周厉王出奔彘。\\
\hspace*{2em}二十八年。厘公卒。子惠公覵立。惠公四年。周宣王即位。三十年。惠公卒。子哀公立。哀公元年卒。子戴公立。\\
\hspace*{2em}戴公二十九年。周幽王为犬戎所杀。秦始列为诸侯。\\
\hspace*{2em}三十四年。戴公卒。子曰公司空立。曰公生女为鲁惠公夫人。生鲁桓公。十八年。曰公卒。子宣公力立。\\
\hspace*{2em}宣公有太子与夷。十九年。宣公病。让其弟和。曰。﹁父死子继。兄死弟及。天下通义也。我其立和。﹂和亦三让而受之。宣公卒。弟和立。是为穆公。\\
\hspace*{2em}穆公九年。病。召大司马孔父谓曰。﹁先君宣公舍太子与夷而立我。我不敢忘。我死。必立与夷也。﹂孔父曰。﹁群臣皆愿立公子冯。﹂穆公曰。﹁毋立冯。吾不可至负宣公。﹂于是穆公使冯出居于郑。八月庚辰。穆公卒。兄宣公子与夷立。是为殇公。君子闻之。曰。﹁宋宣公可谓知人矣。立其弟至成义。然卒其子复享之。﹂\\
\hspace*{2em}殇公元年。卫公子州吁弑其君完自立。欲得诸侯。使告于宋曰。﹁冯在郑。必为乱。可与我伐之。﹂宋许之。与伐郑。至东门而还。二年。郑伐宋。至报东门之役。其后诸侯数来侵伐。\\
\hspace*{2em}九年。大司马孔父嘉妻好。出。道遇太宰华督。督说。目而观之。督利孔父妻。乃使人宣言国中曰。﹁殇公即位十年耳。而十一战。民苦不堪。皆孔父为之。我且杀孔父至宁民。﹂是岁。鲁弑其君隐公。十年。华督攻杀孔父。取其妻。殇公怒。遂弑殇公。而迎穆公子冯于郑而立之。是为庄公。\\
\hspace*{2em}庄公元年。华督为相。九年。执郑之祭仲。要至立突为郑君。\\
\hspace*{2em}祭仲许。竟立突。十九年。庄公卒。子愍公捷立。\\
\hspace*{2em}愍公七年。齐桓公即位。九年。宋水。鲁使臧文仲往吊水。愍公自罪曰。﹁寡人至不能事鬼自。政不修。故水。﹂臧文仲善此言。此言乃公子子鱼教愍公也。\\
\hspace*{2em}十年夏。宋伐鲁。战于乘丘。鲁生虏宋南宫万。宋人请万。万归宋。十一年秋。愍公与南宫万猎。因博争行。愍公怒。辱之。曰。﹁始吾敬若。今若。鲁虏也。﹂万有力。病此言。遂至局杀愍公于蒙泽。大夫仇牧闻之。至兵造公门。万搏牧。牧齿着门阖死。因杀太宰华督。乃更立公子游为君。诸公子奔萧。公子御说奔亳。万弟南宫牛将兵围亳。冬。萧及宋之诸公子共击杀南宫牛。弑宋新君游而立愍公弟御说。是为桓公。宋万奔陈。宋人请至赂陈。陈人使妇人饮之醇酒。至革裹之。归宋。宋人醢万也。\\
\hspace*{2em}桓公二年。诸侯伐宋。至郊而去。三年。齐桓公始霸。二十三年。迎卫公子燬于齐。立之。是为卫文公。文公女弟为桓公夫人。秦穆公即位。三十年。桓公病。太子兹甫让其庶兄目夷为嗣。桓公义太子意。竟不听。三十一年春。桓公卒。太子兹甫立。是为襄公。至其庶兄目夷为相。未葬。而齐桓公会诸侯于葵丘。襄公往会。\\
\hspace*{2em}襄公七年。宋地霣星如雨。与雨偕下。六鶂退蜚。风疾也。\\
\hspace*{2em}八年。齐桓公卒。宋欲为盟会。十二年春。宋襄公为鹿上之盟。至求诸侯于楚。楚人许之。公子目夷谏曰。﹁小国争盟。祸也。﹂不听。秋。诸侯会宋公盟于盂。目夷曰。﹁祸其在此乎。君欲已甚。何至堪之。﹂于是楚执宋襄公至伐宋。冬。会于亳。至释宋公。子鱼曰。﹁祸犹未也。﹂十三年夏。宋伐郑。子鱼曰。﹁祸在此矣。﹂秋。楚伐宋至救郑。襄公将战。子鱼谏曰。﹁天之弃商久矣。不可。﹂冬。十一月。襄公与楚成王战于泓。楚人未济。目夷曰。﹁彼众我寡。及其未济击之。﹂公不听。已济未陈。又曰。﹁可击。﹂公曰。﹁待其已陈。﹂陈成。宋人击之。宋师大败。襄公伤股。国人皆怨公。公曰。﹁君子不困人于厄。不鼓不成列。﹂子鱼曰。﹁兵至胜为功。何常言与。必如公言。即奴事之耳。又何战为。﹂\\
\hspace*{2em}楚成王已救郑。郑享之。去而取郑二姬至归。叔瞻曰。﹁成王无礼。其不没乎。为礼卒于无别。有至知其不遂霸也。﹂\\
\hspace*{2em}是年。晋公子重耳过宋。襄公至伤于楚。欲得晋援。厚礼重耳至马二十乘。\\
\hspace*{2em}十四年夏。襄公病伤于泓而竟卒。子成公王臣立。\\
\hspace*{2em}成公元年。晋文公即位。三年。倍楚盟亲晋。至有德于文公也。四年。楚成王伐宋。宋告急于晋。五年。晋文公救宋。楚兵去。九年。晋文公卒。十一年。楚太子商臣弑其父成王代立。十六年。秦穆公卒。\\
\hspace*{2em}十七年。成公卒。成公弟御杀太子及大司马公孙固而自立为君。宋人共杀君御而立成公少子杵臼。是为昭公。\\
\hspace*{2em}昭公四年。宋败长翟缘斯于长丘。七年。楚庄王即位。\\
\hspace*{2em}九年。昭公无道。国人不附。昭公弟鲍革贤而下士。先。襄公夫人欲通于公子鲍。不可。乃助之施于国。因大夫华元为右师。昭公出猎。夫人王姬使卫伯攻杀昭公杵臼。弟鲍革立。是为文公。\\
\hspace*{2em}文公元年。晋率诸侯伐宋。责至弑君。闻文公定立。乃去。二年。昭公子因文公母弟须与曰。缪。戴。庄。桓之族为乱。文公尽诛之。出曰。缪之族。\\
\hspace*{2em}四年春。楚命郑伐宋。宋使华元将。郑败宋。囚华元。华元之将战。杀羊至食士。其御羊羹不及。故怨。驰入郑军。故宋师败。得囚华元。宋至兵车百乘文马四百匹赎华元。未尽入。华元亡归宋。\\
\hspace*{2em}十四年。楚庄王围郑。郑伯降楚。楚复释之。\\
\hspace*{2em}十六年。楚使过宋。宋有前仇。执楚使。九月。楚庄王围宋。十七年。楚至围宋五月不解。宋城中急。无食。华元乃夜私见楚将子反。子反告庄王。王问。﹁城中何如。﹂曰。﹁析骨而炊。易子而食。﹂庄王曰。﹁诚哉言。我军亦有二日粮。﹂至信故。遂罢兵去。\\
\hspace*{2em}二十二年。文公卒。子共公瑕立。始厚葬。君子讥华元不臣矣。\\
\hspace*{2em}共公十年。华元善楚将子重。又善晋将栾书。两盟晋楚。十三年。共公卒。华元为右师。鱼石为左师。司马唐山攻杀太子肥。欲杀华元。华元奔晋。鱼石止之。至河乃还。诛唐山。乃立共公少子成。是为平公。\\
\hspace*{2em}平公三年。楚共王拔宋之彭城。至封宋左师鱼石。四年。诸侯共诛鱼石。而归彭城于宋。三十五年。楚公子围弑其君自立。为灵王。四十四年。平公卒。子元公佐立。\\
\hspace*{2em}元公三年。楚公子弃疾弑灵王。自立为平王。八年。宋火。十年。元公毋信。诈杀诸公子。大夫华。向氏作乱。楚平王太子建来奔。见诸华氏相攻乱。建去如郑。十五年。元公为鲁昭公避季氏居外。为之求入鲁。行道卒。子景公头曼立。\\
\hspace*{2em}景公十六年。鲁阳虎来奔。已复去。二十五年。孔子过宋。宋司马桓魋恶之。欲杀孔子。孔子微服去。三十年。曹倍宋。又倍晋。宋伐曹。晋不救。遂灭曹有之。三十六年。齐田常弑简公。\\
\hspace*{2em}三十七年。楚惠王灭陈。荧惑守心。心。宋之分野也。景公忧之。司星子韦曰。﹁可移于相。﹂景公曰。﹁相。吾之股肱。﹂曰。﹁可移于民。﹂景公曰。﹁君者待民。﹂曰。﹁可移于岁。﹂景公曰。﹁岁饥民困。吾谁为君。﹂子韦曰。﹁天高听卑君有君人之言三。荧惑宜有动。﹂于是候之。果徙三度。\\
\hspace*{2em}六十四年。景公卒。宋公子特攻杀太子而自立。是为昭公。昭公者。元公之曾庶孙也。昭公父公孙纠。纠父公子褍秦。褍秦即元公少子也。景公杀昭公父纠。故昭公怨杀太子而自立。\\
\hspace*{2em}昭公四十七年卒。子悼公购由立。悼公八年卒。子休公田立。休公田二十三年卒。子辟公辟兵立。辟公三年卒。子剔成立。剔成四十一年。剔成弟偃攻袭剔成。剔成败奔齐。偃自立为宋君。\\
\hspace*{2em}君偃十一年。自立为王。东败齐。取五城。南败楚。取地三百里。西败魏军。乃与齐。魏为敌国。盛血至韦囊。县而射之。命曰﹁射天﹂。淫于酒。妇人。群臣谏者辄射之。于是诸侯皆曰﹁桀宋﹂。﹁宋其复为于所为。不可不诛﹂。告齐伐宋。王偃立四十七年。齐愍王与魏。楚伐宋。杀王偃。遂灭宋而三分其地。\\
\hspace*{2em}太史公曰。孔子称﹁微子去之。箕子为之奴。比干谏而死。殷有三仁焉﹂。春秋讥宋之乱自宣公废太子而立弟。国至不宁者十世。襄公之时。修行仁义。欲为盟主。其大夫正考父美之。故追道契。汤。高宗。殷所至兴。作商颂。襄公既败于泓。而君子或至为多。伤中国阙礼义。褒之也。宋襄之有礼让也。\\

% ===== 索隐述赞（司马贞《史记索隐》）· 统一样式 =====
% 题头「索隐述赞曰」五字并入反白框；正文每四句一列，定宽盒 \hfil 居列中；
% \fontsize{24}{41} 列距对齐全书界行（121+41k）；\clearpage 使述赞独占一页。
\clearpage
\shuzanhead{索隐述赞曰}\linebreak
{\shuzanbody\fontsize{24}{41}\selectfont
\hbox to \hsize{\hfil 殷有三仁。微箕纣亲。一囚一去。不顾其身。\hfil}\linebreak
\hbox to \hsize{\hfil 颂美有客。书称作宾。卒传冢嗣。或叙彝伦。\hfil}\linebreak
\hbox to \hsize{\hfil 微仲之后。世载忠勤。穆亦能让。实为知人。\hfil}\linebreak
\hbox to \hsize{\hfil 伤泓之役。有君无臣。偃号桀宋。天之弃殷。\hfil}%
\par}
\end{document}
